\section{Intuitionistic Proof Transformation}

\subsection{Induction principle for path relation}

\begin{Prop}
$\prove_\HA \Ind(>_\Pi)$.
\end{Prop}

{\em Proof.}

Assume a cyclic proof $\Pi$ is given.

We define a program frame by:
$\Functionnames$ is the set of companion numbers,
and
$\Callsitenames$ is the set of paths from a companion to a bud,
and
$\arity(c)$ for a companion number $c$ is the number of
inductive atomic formulas in the antecedent of the judgment
of companion number $c$.

We define a SC call graph $G$ by:

$V$ is the set of companion numbers in $\Pi$,

$E$ is the set of $(c_1,(S,P),\pi,c_2)$ where
there is a path $\pi$ from a companion of companion number $c_1$ to a bud
of companion number $c_2$ and
\[
S = \{ (i,j) \ |\ F_1(i,j) \}, \\
P = \{ (i,j) \ |\ F_2(i,j) \},
\]
where $F_1(i,j)$ iff
there is some non-progressing trace from
the $i$-th inductive atomic formula in the bottom sequent of $\pi$
to the $j$-the inductive atomic formula in the top sequent of $\pi$,
and
$F_2(i,j)$ iff
there is non-progressing trace from
the $i$-th inductive atomic formula in the bottom sequent of $\pi$
to the $j$-the inductive atomic formula in the top sequent of $\pi$.

Since the cyclic proof $\Pi$ satisfies, the SC call graph $G$ satisfies
SCT condition.

Claim 1. $\<c_1,\<\Vec x\>\> >_\Pi \<c_2,\<\Vec y\>\>$ implies
$\rho(c_1,\Vec x) >_\Lex \rho(c_2,\Vec y)$.

We can show it as follows.
From $\<c_1,\<\Vec x\>\> >_\Pi \<c_2,\<\Vec y\>\>$,
there is a path $\pi$ from the bottom sequent of companion number $c_1$
to the top sequent of companion number $c_2$ and
$\<c_1,\<\Vec x\>\> >_\pi \<c_2,\<\Vec y\>\>$.
Then there is a SC graph $R=(c_1,-,\pi,c_2)$ in the SC call graph $G$
and $\Vec x R \Vec y$.
Hence $(c_1,\Vec x) \to^\pi (c_2,\Vec y)$.
By Theorem \ref{th:rankingfunc}, we have
$\rho(f, \Vec x) >_\Lex (g, \Vec y)$.
We have shown the claim 1.

We have the induction principle for lexicographic ordering on natural number
\[
\forall \Vec xf((\forall \Vec yg(\rho(g,\Vec y) <_\Lex \rho(f,\Vec x) \imp A(g,\Vec y)))
\imp A(f,\Vec x))
\imp \forall \Vec xf A(f,\Vec x).
\]
By clamin 1, taking $x$ to be $\<\Vec x\>$,
it implies the induction principle for the path relation
\[
\forall xf((\forall yg(\<g,y\> <_\Pi \<f,x\> \imp A(\<g,y\>)))
\imp A(\<f,x\>))
\imp \forall xf A(\<f,x\>).
\]
$\Box$
